\documentclass[12pt]{article}
\usepackage[margin=2cm]{geometry}
\setlength{\parindent}{1em}
\setlength{\parskip}{1em}

\usepackage{amsmath}
\usepackage{booktabs,dcolumn,setspace,float,multirow}

\begin{document}

\section{Interpretation and Purpose of Sufficient Statistics}

\subsection{Markov Transition Statistics}

\paragraph{$C_1$ and $C_0$ (Initial State Counts)}

\textbf{What they are:}
\begin{itemize}
\item $C_1$: Expected weighted count of individuals employed at wave 1 ($h_1 = 1$)
\item $C_0$: Expected weighted count of individuals nonemployed at wave 1 ($h_1 = 0$)
\end{itemize}

\textbf{Formula:}
\begin{align*}
C_1 &= \sum_{i,h} w_i\gamma_{ih}\cdot h_1,\\
C_0 &= \sum_{i,h} w_i\gamma_{ih}\cdot(1-h_1).
\end{align*}

\textbf{Why calculated:}
\begin{itemize}
\item \emph{Purpose}: Estimate initial employment probability $\alpha = \Pr(h_1 = 1)$
\item \emph{M-step formula}: $\hat{\alpha} = C_1 / (C_1 + C_0)$
\item \emph{Interpretation}: ``Of all people at wave 1, what fraction were employed?''
\end{itemize}

\textbf{Example:}
If $C_1 = 600$ and $C_0 = 400$ (total survey weight = 1000):
\[
\hat{\alpha} = 600/1000 = 0.60 \quad\text{(60\% initially employed)}.
\]

\paragraph{$D_1$ and $D_0$ (Transition Exposure/Denominators)}

\textbf{What they are:}
\begin{itemize}
\item $D_1$: Expected weighted count of ``employment episodes'' (times when someone was employed at $t-1$)
\item $D_0$: Expected weighted count of ``nonemployment episodes'' (times when someone was nonemployed at $t-1$)
\end{itemize}

\textbf{Formula:}
\begin{align*}
D_1 &= \sum_{i,h} w_i\gamma_{ih}\sum_{t=2}^3\mathbf{1}\{h_{t-1}=1\},\\
D_0 &= \sum_{i,h} w_i\gamma_{ih}\sum_{t=2}^3\mathbf{1}\{h_{t-1}=0\}.
\end{align*}

\textbf{Why calculated:}
\begin{itemize}
\item \emph{Purpose}: Denominator for transition probabilities (exposure time)
\item \emph{M-step formulas}: 
  \begin{align*}
  \hat{\theta}_1 &= T_{11} / D_1 \quad\text{(stay-employed rate)},\\
  \hat{\theta}_0 &= T_{01} / D_0 \quad\text{(job-finding rate)}.
  \end{align*}
\item \emph{Interpretation}: ``How many times were people `at risk' of transitioning from each state?''
\end{itemize}

\textbf{Example:}
Person with $h = (1,1,0)$:
\begin{itemize}
\item At $t=1$: employed (contributes 1 to $D_1$)
\item At $t=2$: employed (contributes 1 to $D_1$)
\item Total contribution to $D_1$: 2
\end{itemize}
Across all $N=500$ people: $D_1 = 850$ $\Rightarrow$ There were 850 employment episodes where transitions could occur.

\textbf{Note:} We only use $t=1,2$ (not $t=3$) because there is no $t=4$ to transition to.

\paragraph{$T_{11}$ (Stay-Employed Transitions)}

\textbf{What it is:}
\begin{itemize}
\item Expected weighted count of transitions where someone \textbf{stayed employed} ($h_{t-1} = 1 \to h_t = 1$)
\end{itemize}

\textbf{Formula:}
\[
T_{11} = \sum_{i,h} w_i\gamma_{ih}\sum_{t=2}^3\mathbf{1}\{h_{t-1}=1, h_t=1\}.
\]

\textbf{Why calculated:}
\begin{itemize}
\item \emph{Purpose}: Numerator for employment persistence probability
\item \emph{M-step formula}: $\hat{\theta}_1 = T_{11} / D_1$
\item \emph{Interpretation}: ``Of all employment episodes, how many resulted in staying employed?''
\end{itemize}

\textbf{Example:}
\[
T_{11} = 750,\, D_1 = 850 \quad\Rightarrow\quad \hat{\theta}_1 = 750/850 = 0.882 \quad\text{(88.2\% of employed stay employed next quarter)}.
\]

\paragraph{$T_{01}$ (Job-Finding Transitions)}

\textbf{What it is:}
\begin{itemize}
\item Expected weighted count of transitions where someone \textbf{found employment} ($h_{t-1} = 0 \to h_t = 1$)
\end{itemize}

\textbf{Formula:}
\[
T_{01} = \sum_{i,h} w_i\gamma_{ih}\sum_{t=2}^3\mathbf{1}\{h_{t-1}=0, h_t=1\}.
\]

\textbf{Why calculated:}
\begin{itemize}
\item \emph{Purpose}: Numerator for job-finding probability
\item \emph{M-step formula}: $\hat{\theta}_0 = T_{01} / D_0$
\item \emph{Interpretation}: ``Of all nonemployment episodes, how many resulted in finding a job?''
\end{itemize}

\textbf{Example:}
\[
T_{01} = 120,\, D_0 = 650 \quad\Rightarrow\quad \hat{\theta}_0 = 120/650 = 0.185 \quad\text{(18.5\% of nonemployed find jobs next quarter)}.
\]

\paragraph{$M$ (Total Misclassifications)}

\textbf{What it is:}
\begin{itemize}
\item Expected weighted count of \textbf{all state misclassifications} across all waves
\end{itemize}

\textbf{Formula:}
\[
M = \sum_{i,h} w_i\gamma_{ih}\sum_{t=1}^3 \mathbf{1}\{s_{it} \neq h_t\}.
\]

\textbf{Why calculated:}
\begin{itemize}
\item \emph{Purpose}: Estimate misclassification probability $\pi$
\item \emph{M-step formula}: $\hat{\pi} = M / (3\cdot\sum_i w_i)$ (divide by $3N$ because each person has 3 waves)
\item \emph{Interpretation}: ``What fraction of all state reports (across all people and waves) were errors?''
\end{itemize}

\textbf{Example:}
\[
M = 150,\, \text{total survey weight} = 1000,\, \text{number of wave-reports} = 3000 \quad\Rightarrow\quad \hat{\pi} = 150/3000 = 0.05 \quad\text{(5\% error rate)}.
\]

\subsection{Duration Statistics}

\paragraph{$S_g$ and $N_g$ (Employment Continuation Variance)}

\textbf{What they are:}
\begin{itemize}
\item $S_g$: Sum of squared \textbf{employment tenure increments} (for continuations only)
\item $N_g$: Count of employment continuations (number of increment observations)
\end{itemize}

\textbf{Formula:}
\begin{align*}
S_g &= \sum_{i,h} w_i\gamma_{ih}\sum_{t=2}^3 \mathbf{1}\{h_{t-1}=1, h_t=1, s_{it}=1, s_{i,t-1}=1\} \cdot \big[\Delta^g_{it}\big]^2,\\
N_g &= \sum_{i,h} w_i\gamma_{ih}\sum_{t=2}^3 \mathbf{1}\{h_{t-1}=1, h_t=1, s_{it}=1, s_{i,t-1}=1\},
\end{align*}
where $\Delta^g_{it} = (g_{it} - g_{i,t-1}) - 0.25$ from \eqref{eq:emp_increment}. The condition $s_{i,t-1}=1$ ensures the lagged tenure $g_{i,t-1}$ is observed, so the increment is well-defined. When the previous wave is misclassified ($s_{i,t-1}=0$), the increment is undefined and that observation is excluded.

\textbf{Why calculated:}
\begin{itemize}
\item \emph{Purpose}: Estimate employment measurement error variance $\sigma_g^2$
\item \emph{M-step formula}: $\hat{\sigma}_g^2 = (S_g + 2S_g^{\mathrm{start}}) / (2(N_g + N_g^{\mathrm{start}}))$ \eqref{eq:mstep_sigmas} \quad$\leftarrow$ \textbf{Pools continuations and starts}
\item \emph{Interpretation}: ``How noisy are reported tenure increments?''
\end{itemize}

\textbf{Example:}
If $S_g = 0.80$, $N_g = 200$, $S_g^{\mathrm{start}} = 0.05$, $N_g^{\mathrm{start}} = 30$:
\[
\hat{\sigma}_g^2 = \frac{0.80 + 2(0.05)}{2(200+30)} = \frac{0.90}{460} = 0.00196 \quad\Rightarrow\quad \hat{\sigma}_g \approx 0.044\text{ years} \approx 16\text{ days}.
\]
When $N_g^{\mathrm{start}}=0$, this reduces to $\hat{\sigma}_g^2 = S_g/(2N_g)$.

\textbf{Why the factor of 2?}

The increment $\Delta^g_{it} = (g_{it} - g_{i,t-1}) - 0.25$ has \textbf{two independent measurement errors}:
\begin{align*}
\Delta^g_{it} &= \big[g_t^* + \varepsilon^g_{it}\big] - \big[g_{t-1}^* + \varepsilon^g_{i,t-1}\big] - 0.25\\
&= 0.25 + \varepsilon^g_{it} - \varepsilon^g_{i,t-1} - 0.25\\
&= \varepsilon^g_{it} - \varepsilon^g_{i,t-1}.
\end{align*}
Therefore,
\[
\mathrm{Var}(\Delta^g_{it}) = \mathrm{Var}(\varepsilon^g_{it}) + \mathrm{Var}(\varepsilon^g_{i,t-1}) = \sigma_g^2 + \sigma_g^2 = 2\sigma_g^2.
\]
So: $\mathbb{E}\big[(\Delta^g_{it})^2\big] = 2\sigma_g^2$, and solving gives $\sigma_g^2 = \mathbb{E}\big[(\Delta^g_{it})^2\big] / 2 = S_g / (2\cdot N_g)$.

\paragraph{$S_d$ and $N_d$ (Nonemployment Continuation Variance)}

\textbf{What they are:}
\begin{itemize}
\item $S_d$: Sum of squared \textbf{nonemployment duration increments}
\item $N_d$: Count of nonemployment continuations
\end{itemize}

\textbf{Formula:}
\begin{align*}
S_d &= \sum_{i,h} w_i\gamma_{ih}\sum_{t=2}^3 \mathbf{1}\{h_{t-1}=0, h_t=0, s_{it}=0, s_{i,t-1}=0\} \cdot \big[\Delta^d_{it}\big]^2,\\
N_d &= \sum_{i,h} w_i\gamma_{ih}\sum_{t=2}^3 \mathbf{1}\{h_{t-1}=0, h_t=0, s_{it}=0, s_{i,t-1}=0\},
\end{align*}
where $\Delta^d_{it} = (d_{it} - d_{i,t-1}) - 0.25$ from \eqref{eq:nonemp_increment}. The condition $s_{i,t-1}=0$ ensures the lagged nonemployment duration $d_{i,t-1}$ is observed.

\textbf{Why calculated:}
\begin{itemize}
\item \emph{Purpose}: Estimate nonemployment measurement error variance $\sigma_d^2$
\item \emph{M-step formula}: $\hat{\sigma}_d^2 = (S_d + 2S_d^{\mathrm{start}}) / (2(N_d + N_d^{\mathrm{start}}))$ \eqref{eq:mstep_sigmas}
\item \emph{Interpretation}: ``How noisy are reported nonemployment duration increments?''
\end{itemize}

Same logic as employment, but for the nonemployment regime.

\paragraph{$S_g^{\mathrm{start}}$, $N_g^{\mathrm{start}}$, $S_d^{\mathrm{start}}$, $N_d^{\mathrm{start}}$ (Within-Panel Start Statistics)}

\textbf{What they are:}
\begin{itemize}
\item $S_g^{\mathrm{start}}$: Sum of squared deviations $(g_{it}-0.25)^2$ at matched employment starts
\item $N_g^{\mathrm{start}}$: Count of matched employment starts
\item $S_d^{\mathrm{start}}, N_d^{\mathrm{start}}$: Analogous for nonemployment starts
\end{itemize}

\textbf{Formulas:}

\emph{Employment starts} ($t\ge 2$, $h_{t-1}=0$, $h_t=1$, $s_{it}=1$):
\begin{align*}
S_g^{\mathrm{start}} &= \sum_{i,h} w_i\gamma_{ih}\sum_{t=2}^3\mathbf{1}\{h_{t-1}=0,\,h_t=1,\,s_{it}=1\}\big(g_{it}-0.25\big)^2,\\
N_g^{\mathrm{start}} &= \sum_{i,h} w_i\gamma_{ih}\sum_{t=2}^3\mathbf{1}\{h_{t-1}=0,\,h_t=1,\,s_{it}=1\}.
\end{align*}

\emph{Nonemployment starts} ($t\ge 2$, $h_{t-1}=1$, $h_t=0$, $s_{it}=0$):
\begin{align*}
S_d^{\mathrm{start}} &= \sum_{i,h} w_i\gamma_{ih}\sum_{t=2}^3\mathbf{1}\{h_{t-1}=1,\,h_t=0,\,s_{it}=0\}\big(d_{it}-0.25\big)^2,\\
N_d^{\mathrm{start}} &= \sum_{i,h} w_i\gamma_{ih}\sum_{t=2}^3\mathbf{1}\{h_{t-1}=1,\,h_t=0,\,s_{it}=0\}.
\end{align*}

\textbf{Why calculated:}
\begin{itemize}
\item \emph{Purpose}: Provide additional information about $\sigma_g^2$ and $\sigma_d^2$, pooled with continuation statistics in the M-step
\item \emph{Distribution}: At a within-panel start, the true duration is exactly $0.25$ (one full quarter). The observed duration is $(0.25 + \varepsilon)$ where $\varepsilon \sim N(0,\sigma^2)$, so $(g_{it}-0.25) \sim N(0, \sigma_g^2)$. This is a standard normal-variance sufficient statistic.
\item \emph{Pooled MLE}: The factor of 2 on $S^{\mathrm{start}}$ in \eqref{eq:mstep_sigmas} rescales start contributions (variance $\sigma^2$) to the continuation scale (variance $2\sigma^2$). Derivation: the joint profile log-likelihood is
\[
\ell(v) = -\tfrac{N_g}{2}\log(2v) - \frac{S_g}{4v} - \tfrac{N_g^{\mathrm{start}}}{2}\log(v) - \frac{S_g^{\mathrm{start}}}{2v},
\]
and setting $\partial\ell/\partial v=0$ gives $\hat{\sigma}_g^2 = (S_g + 2S_g^{\mathrm{start}})/(2(N_g + N_g^{\mathrm{start}}))$.
\item \emph{No $s_{i,t-1}$ condition needed}: Unlike continuations, starts do not use a lagged duration, so there is no requirement on the previous wave's observed state.
\end{itemize}

\textbf{Example:}
If $S_g^{\mathrm{start}} = 0.05$, $N_g^{\mathrm{start}} = 30$, combined with $S_g = 0.80$, $N_g = 200$:
\[
\hat{\sigma}_g^2 = \frac{0.80 + 2(0.05)}{2(200+30)} = \frac{0.90}{460} = 0.00196.
\]

\paragraph{Exponential rates $\lambda_g, \lambda_d$ (Linked CTMC)}

\textbf{What they are:}
\begin{itemize}
\item Under the linked specification, $\lambda_g$ and $\lambda_d$ are \textbf{not estimated from separate sufficient statistics}
\item They are determined entirely by the transition probability estimates via the CTMC link \eqref{eq:lambdas_link}
\end{itemize}

\textbf{M-step formula:}
\[
\hat{\lambda}_g = -\tfrac{1}{3}\log(1-\hat{\theta}_1), \qquad \hat{\lambda}_d = -\tfrac{1}{3}\log(1-\hat{\theta}_0).
\]

\textbf{Why no separate statistics?}
\begin{itemize}
\item Misclassified durations and wave-1 durations follow the EMG (exGaussian) distribution, not pure Exponential
\item The EMG distribution has no closed-form MLE for the exponential rate parameter
\item Under the linked specification, $\lambda$ is pinned down by $\theta$, so no numerical optimization is needed
\item These EMG observations still influence the E-step by entering the emission densities $\ell_t(\cdot)$ and thereby shaping the responsibilities $\gamma_{ih}$
\end{itemize}

\textbf{Example:}
\[
\hat{\theta}_1 = 0.882 \quad\Rightarrow\quad \hat{\lambda}_g = -\tfrac{1}{3}\log(1-0.882) = -\tfrac{1}{3}\log(0.118) = 0.713.
\]
This implies a mean true employment spell length of $1/\hat{\lambda}_g = 1.40$ years $\approx$ 17 months.

\subsection{Summary Table}

\begin{table}[h]
\centering
\small
\begin{tabular}{llll}
\hline
\textbf{Statistic} & \textbf{What It Counts} & \textbf{M-Step Use} & \textbf{Estimate} \\
\hline
$C_1, C_0$ & Initial employment/nonemployment & Initial probability & $\alpha = C_1/(C_1+C_0)$ \\
$D_1, D_0$ & Employment/nonemployment exposure & Transition denominators & (used below) \\
$T_{11}$ & Stay-employed transitions & Persistence numerator & $\theta_1 = T_{11}/D_1$ \\
$T_{01}$ & Job-finding transitions & Job-finding numerator & $\theta_0 = T_{01}/D_0$ \\
$M$ & Total misclassifications & Error rate & $\pi = M/(3\sum_i w_i)$ \\
$S_g, N_g$ & Tenure continuation increments & \multirow{2}{*}{Pooled $\to$ $\sigma_g^2$} & \multirow{2}{*}{$\frac{S_g+2S_g^{\mathrm{start}}}{2(N_g+N_g^{\mathrm{start}})}$} \\
$S_g^{\mathrm{start}}, N_g^{\mathrm{start}}$ & Employment start deviations & & \\
$S_d, N_d$ & Nonemp continuation increments & \multirow{2}{*}{Pooled $\to$ $\sigma_d^2$} & \multirow{2}{*}{$\frac{S_d+2S_d^{\mathrm{start}}}{2(N_d+N_d^{\mathrm{start}})}$} \\
$S_d^{\mathrm{start}}, N_d^{\mathrm{start}}$ & Nonemp start deviations & & \\
(none) & $\lambda_g,\lambda_d$ via CTMC link & Deterministic & $\lambda_g = -\tfrac{1}{3}\log(1-\theta_1)$ \\
\hline
\end{tabular}
\end{table}

\newpage
\subsection{Key Conceptual Points}

\begin{enumerate}
\item \textbf{``Expected'' counts}: All statistics involve $\gamma_{ih}$ (posterior probability), so they are \textbf{weighted by uncertainty} about the true history.

\item \textbf{Survey weights}: All statistics multiply by $w_i$ to account for sampling design.

\item \textbf{Two sources for $\sigma^2$}: Continuation increments ($S_g, N_g$) and within-panel starts ($S_g^{\mathrm{start}}, N_g^{\mathrm{start}}$) both identify $\sigma_g^2$ and are pooled in the M-step \eqref{eq:mstep_sigmas}. Continuations require $s_{i,t-1}=s_{it}$ so that both current and lagged durations are observed; when the previous wave is misclassified, the increment is undefined and that observation is excluded. Starts do not need this condition because they use the level $(g_{it}-0.25)$ rather than a lagged difference.

\item \textbf{No separate statistics for $\lambda$}: Under the linked specification, $\lambda_g$ and $\lambda_d$ are determined by $\theta_1$ and $\theta_0$ via the CTMC link \eqref{eq:lambdas_link}. Misclassified and wave-1 durations follow the EMG distribution, which has no closed-form MLE for the rate parameter, so no separate sufficient statistics are needed or computed.

\item \textbf{EMG enters through responsibilities}: Wave-1 EMG observations, continuation-with-previous-misclassified cases, and misclassified durations all enter the emission density $\ell_t(\cdot)$ and thereby influence the responsibilities $\gamma_{ih}$, but they do not yield closed-form sufficient statistics for $\lambda$. Their information shapes the soft assignments, which in turn affect all other sufficient statistics.
\end{enumerate}

\end{document}